\documentclass{article}
\usepackage[utf8]{inputenc}
\usepackage[T1]{fontenc}
\usepackage{geometry}
\usepackage{graphicx}
\usepackage{float}
\usepackage{hyperref}
\usepackage[english,indonesian]{babel}

\title{Optimasi Pemrosesan Template LaTeX Menggunakan Pendekatan Berbasis AI}
\author{Academic Support Agent}
\date{\today}

\begin{document}
\maketitle

\begin{abstract}

\end{abstract}

\chapter{PENDAHULUAN}


\chapter{TINJAUAN PUSTAKA}


% Original content (commented out):
% \documentclass{article}
\usepackage[utf8]{inputenc}

\usepackage[ngerman,english]{babel}
\usepackage{fontawesome}
\usepackage{biblatex}
\bibliography{references.bib}
\usepackage{hyperref}
\usepackage{graphicx}

\usepackage{amsmath}
\usepackage{gensymb} % for \degreee command
\usepackage{csquotes}

\usepackage{tikz}

\usepackage{xcolor}
\definecolor{alert}{HTML}{CD5C5C}
\definecolor{myblue}{HTML}{616AC5}


\usepackage{tcolorbox}
\tcbuselibrary{most}

 \tcbset{tcbox width=auto,left=1mm,top=1mm,bottom=1mm,
right=1mm,boxsep=1mm,middle=1pt}

\newenvironment{mycolorbox}[2]{%
\begin{tcolorbox}[capture=minipage,fonttitle=\normalsize, enhanced jigsaw,boxsep=1mm,colback=#1!50!white,on line,tcbox width=auto, toptitle=0mm,colframe=#1,opacityback=0.7,nobeforeafter,title=#2]%
}{\end{tcolorbox}}


%--------------------------------

\title{Fast typing with LaTeX}
\author{Marie Musterfrau}
\date{\today}

\begin{document}


\maketitle
\tableofcontents

\section{Introduction}
This is a test sentence.

This is another one.

\subsection{test}
Powered by \href{http://latex-ninja.com}{the \LaTeX{} Ninja}.

\textit{makes text italic}, a similar but different effect is achieved with \emph{which highlights}. Make some \textbf{bold} statements and don't forget the ``quotes''.\newpage will cause a manual pagebreak. Now copy this sentence into your test document.

\begin{verbatim}
@Article{CiteKey,
author = {Max Mustermann and Marie Musterfrau},
title = {Important Findings on Fast Typing in \LaTeX{}},
journal = {The LaTeX Ninja Blog},
year = {2019},
volume = {1},
number = {3},
pages = {22--23},
url = {https://latex-ninja.com}
}
\end{verbatim}
% What a useless example. Great I can hate on people in LaTeX using comments and they'll never know.

\verb|\emph{}| is used for highlighting.

\texttt{some code} always looks cooler in a typewriter font.
\textsc{Authors} really love it when their names are set in small caps.


\emph{I just wanted to add that somestimes people have stuff \emph{so incredibly important} to say that they want multiple layers of highlighting possible.} Thank God, \textit{unlike textit}, \verb|\emph{}| is totally up for the job.

\begin{itemize}
    \item used to note down some stuff using bullet points $\to$ very practical
    \item additional item
\end{itemize}


\begin{enumerate}
    \item note down stuff in ordered form
    \item use the \verb|enumerate| environment
\end{enumerate}


\newpage
\section{In case you need TikZ}

\begin{figure}
    \centering
    \begin{tikzpicture}
\scriptsize
\node[](Kreismittelpunkt) at (0,0) {};
\node[](K1) at (0,6) {};
\node[](K2) at (4,5) {};

\draw[] (Kreismittelpunkt) -- ++(0:5cm) node[](C){} -- ++(0:5cm) node[](K2){};
\draw[] (Kreismittelpunkt) -- ++(30:4cm)  node[](A){}  --  ++(30:2cm)  node[](B){}  -- ++(30:4cm) node[](K1){};
\draw[] (K2) -- (K1);
\draw[fill=myblue] (K2) circle (2.5pt) node[below=0.5em]{K\textsubscript{2}};
\draw[fill=myblue] (K1) circle (2.5pt) node[above=0.5em]{K\textsuperscript{1}};

\draw[fill=myblue] (A) circle (2.5pt) node[above=0.5em]{A};
\draw[fill=myblue] (B) circle (2.5pt) node[above=0.5em]{B};
\draw[fill=myblue] (C) circle (2.5pt) node[below=0.5em]{C};
\draw[] (A) -- (C);
\draw[] (C) -- (B);
\draw[] (B) -- (K2);

% coordinates
\coordinate (K2) at (K2);
\coordinate (K1) at (K1);
\coordinate (A) at (A);
\coordinate (B) at (B);
\coordinate (C) at (C);

\draw[fill=myblue] (Kreismittelpunkt) circle (2.5pt) node[below=0.5em]{K};

\draw[-,draw=black!70] (K2) to[bend right=12] (K1);
\end{tikzpicture}
    \caption{a sketch}
    \label{fig:sketch}
\end{figure}

\normalsize

\clearpage
\section{Tables, graphics, etc.}

\begin{mycolorbox}{myblue}{Satz von der Winkelsumme in der euklidischen Geometrie}
Liegt das Polygon in einer euklidischen Ebene, ist die Winkelsumme durch die Formel
\begin{align*}
(n \cdot 180\degree ) - 360\degree &= (n-2) \cdot 180\degree \\
\end{align*}
gegeben, wobei $n$ für die Zahl der Ecken des Polygons steht.
\end{mycolorbox}
\smallskip

\vspace{1cm}

\includegraphics[width=0.8\textwidth]{test.png}


Like shown in \ref{fig:sketch}, \LaTeX{} is really cool.\footcite[cf.][33--55]{mycitation}

\section{Feel free to add stuff according to your own needs}



\begin{thebibliography}{99}

\bibitem{Principe} Lawrence R. \textbf{Principe}, {\it The Secrets of Alchemy} (Chicago 2013).

\bibitem{Comment} Sometimes, people have extra things to say. 

\end{thebibliography}

\end{document}


\end{document}